\cvsection{Expériences professionnelles}

\begin{cventries}
  \cventry
    {Stagiaire ingénieur de recherche en IA}
    {INRIA}
    {Bordeaux, France}
    {Mars 2026 -- Septembre 2026}
    {
      \begin{cvitems}
        \item {Conception en autonomie, avec le suivi de mes encadrants, d’un agent IA multi-étapes destiné à aider les chercheurs à trouver des articles scientifiques.}
        \item {Construction de l’agent sans framework avec intégration de l’API OpenAI, conception de prompts et gestion des outils, de la mémoire et du contexte.}
        \item {Collecte automatisée de documents scientifiques en ligne, préparation et découpage en segments adaptés au pipeline RAG.}
        \item {Génération des embeddings et indexation des segments dans la base vectorielle ChromaDB.}
        \item {Recherche sémantique des passages pertinents puis intégration du contexte récupéré dans les requêtes adressées au LLM.}
        \item {Développement du backend Python/FastAPI avec API REST, OpenAPI et MongoDB, puis intégration avec Docker et GitHub CI/CD.}
        \item {Installation et configuration de modèles locaux, puis utilisation de vLLM pour les lancer et les servir sur des GPU NVIDIA.}
        \item {Comparaison de Qwen, Gemma et GLM et de plusieurs quantifications sur infrastructure HPC avec Slurm selon la latence, le TTFT, la qualité et la taille du contexte.}
        \item {Définition du protocole expérimental, analyse des résultats et contribution à un article scientifique comparant les modèles.}
      \end{cvitems}
    }

  \cventry
    {Stagiaire de recherche en apprentissage par renforcement}
    {Université Karunya}
    {Coimbatore, Tamil Nadu, Inde}
    {Juin 2025 -- Septembre 2025}
    {
      \begin{cvitems}
        \item {Développement d’un modèle d’apprentissage par renforcement et de son environnement d’entraînement pour la navigation autonome d’un robot hospitalier.}
        \item {Déploiement et optimisation du logiciel et du modèle sur un Raspberry Pi embarqué.}
      \end{cvitems}
    }
\end{cventries}
